% Figure 11.1: one-loop vacuum-polarization diagram.
% Source: physical PDF p. 497 (printed p. 474).
\begingroup
\renewcommand{\thefigure}{11.\arabic{figure}}
\setcounter{figure}{0}
\begin{figure}[tbp]
\centering
\begin{tikzpicture}[
  electron/.style={line width=0.65pt,
    postaction={decorate},decoration={markings,
      mark=at position 0.52 with {\arrow{Stealth[length=4.5pt,width=3.4pt]}}}},
  photon/.style={line width=0.6pt,decorate,
    decoration={snake,amplitude=1.05pt,segment length=5.2pt}},
  every node/.style={font=\small}
]
  \draw[photon] (-4.2,0) -- (-1.55,0);
  \draw[photon] (1.55,0) -- (4.2,0);
  \draw[electron] (-1.55,0) arc[start angle=180,end angle=0,x radius=1.55cm,y radius=1.15cm];
  \draw[electron] (1.55,0) arc[start angle=0,end angle=-180,x radius=1.55cm,y radius=1.15cm];
  \node[above] at (0,1.23) {$p$};
  \node[below] at (0,-1.22) {$p-q$};
  \node at (-2.82,-0.30) {$q$};
  \draw[-{Stealth[length=4pt,width=3pt]},line width=0.55pt] (-3.07,-0.66) -- (-2.57,-0.66);
\end{tikzpicture}
\caption{The one-loop diagram for the vacuum polarization in quantum electrodynamics. Here wavy lines represent photons; lines carrying arrows represent electrons.}
\label{fig:11.1}
\end{figure}
\endgroup
